\documentclass[a4paper,11pt]{article}
\usepackage[english]{babel}
\usepackage[utf8]{inputenc}
\usepackage[T1]{fontenc}
\usepackage[a4paper,margin=1.5cm,top=1.2cm,bottom=1.2cm]{geometry}
\usepackage{graphicx}
\usepackage{amsmath}
\usepackage{booktabs}
\usepackage{siunitx}
\usepackage{float}
\usepackage{microtype}
\usepackage{newtxtext,newtxmath}
\usepackage{setspace}
\usepackage{fancyhdr}
\setstretch{1.15}

\bibliographystyle{apsrev}
\setcounter{secnumdepth}{3}
\pagestyle{fancy}
\fancyhf{}
\fancyfoot[C]{\thepage}
\renewcommand{\headrulewidth}{0pt}
\renewcommand{\footrulewidth}{0pt}
\setlength{\headheight}{14pt}
\setlength{\footskip}{14pt}

\begin{document}

\title{[TITLE]}
\author{[AUTHOR 1] \and [AUTHOR 2]}
\maketitle

\begin{abstract}
[ABSTRACT: 4-6 sentences, no numerical values unless explicitly requested.]
\end{abstract}

\section{Introduction and Objectives}

\subsection{Theoretical Background}
[PLACEHOLDER: concise theoretical framework for the experiment.]

% Example equations block (replace with experiment-specific model)
\begin{equation}
[MODEL EQUATION 1]
\end{equation}

\begin{equation}
[MODEL EQUATION 2]
\end{equation}

\subsection{Measurement Methodology}
[PLACEHOLDER: apparatus operation, measurement modes, and data acquisition workflow.]

\section{Materials and Methods}

\subsection{Experimental Setup}
[PLACEHOLDER: apparatus and sensor description.]

\subsection{Data Acquisition and Analysis}
[PLACEHOLDER: describe data source only (human-provided script arrays or Excel fields).]

\section{Results and Data Analysis}

\subsection{Measured Frequencies}
Table~\ref{tab:frequencies} summarizes the characteristic frequencies.

\begin{table}[H]
\centering
\begin{tabular}{lccc}
\toprule
\textbf{Mode} & \textbf{FFT (\unit{rad/s})} & \textbf{Fit (\unit{rad/s})} & \textbf{Description} \\
\midrule
Natural ($\omega_0$) & [TODO_VALUE $\pm$ TODO_UNCERTAINTY] & [TODO_VALUE $\pm$ TODO_UNCERTAINTY] & Center reference oscillator \\
Symmetric ($\omega_s$) & [TODO_VALUE $\pm$ TODO_UNCERTAINTY] & [TODO_VALUE $\pm$ TODO_UNCERTAINTY] & Both oscillators in phase \\
Asymmetric ($\omega_a$) & [TODO_VALUE $\pm$ TODO_UNCERTAINTY] & [TODO_VALUE $\pm$ TODO_UNCERTAINTY] & Oscillators out of phase \\
\bottomrule
\end{tabular}
\caption{Measured angular frequencies. Values must come from human-provided data with explicit provenance.}
\label{tab:frequencies}
\end{table}

\subsection{Mode Superposition}
\begin{table}[H]
\centering
\begin{tabular}{lccc}
\toprule
\textbf{Mode} & \textbf{FFT (\unit{rad/s})} & \textbf{Fit (\unit{rad/s})} & \textbf{Description} \\
\midrule
Minimum ($\omega_{min}$) & [TODO_VALUE $\pm$ TODO_UNCERTAINTY] & [TODO_VALUE $\pm$ TODO_UNCERTAINTY] & Superposition slow mode \\
Maximum ($\omega_{max}$) & [TODO_VALUE $\pm$ TODO_UNCERTAINTY] & [TODO_VALUE $\pm$ TODO_UNCERTAINTY] & Superposition fast mode \\
\bottomrule
\end{tabular}
\caption{Extracted angular frequencies from superposed motion.}
\label{tab:superposition_frequencies}
\end{table}

% [FIGURE: Insert frequency-response plots and superposition/beating visualization]

\section{Discussion and Conclusions}
% [TODO: Human investigator must complete this section]
% [TODO: If value source is uncertain, keep placeholders and do not infer]

\begin{equation}
\omega_0 = [TODO], \qquad \omega_s = [TODO], \qquad \omega_a = [TODO].
\end{equation}

\begin{equation}
\frac{\omega_s}{\omega_0} = [TODO], \qquad
\frac{\omega_a}{\omega_0} = [TODO], \qquad
\frac{\omega_a}{\omega_s} = [TODO].
\end{equation}

\bibliography{biblio}

\end{document}
